% =====================================================================
%  An Independent Reproduction and Differential Evaluation of
%  GlitchProber (ASE'24) and GlitchCleaner (AAAI-26)
%
%  Self-contained study document. Compile with pdflatex (twice for refs).
%  Every number in this document is produced by the code in this
%  repository and traceable to a file under results/; see Appendix A.
% =====================================================================
\documentclass[11pt,a4paper]{article}

\usepackage[utf8]{inputenc}
\usepackage[T1]{fontenc}
\usepackage{lmodern}
\usepackage[margin=2.5cm]{geometry}
\usepackage{amsmath,amssymb}
\usepackage{booktabs}
\usepackage{longtable}
\usepackage{graphicx}
\usepackage{xcolor}
\usepackage{caption}
\usepackage[hidelinks]{hyperref}
\usepackage{enumitem}
\usepackage{fancyvrb}

\newcommand{\gp}{\textsc{GlitchProber}}
\newcommand{\gc}{\textsc{GlitchCleaner}}
\newcommand{\Ppaper}{$P_{\text{paper}}$}
\newcommand{\Pcode}{$P_{\text{code}}$}

\title{\textbf{An Independent Reproduction and Differential Evaluation of\\
Two Glitch-Token Mitigation Methods}\\[0.3em]
\large \gp{} (ASE'24) and \gc{} (AAAI-26)}
\author{Benchmark study --- methodology, execution, and results}
\date{}

\begin{document}
\maketitle

\begin{abstract}
\noindent
Two published methods claim to detect and repair \emph{glitch tokens} --- vocabulary
entries so poorly trained that they derail a language model's output. \gp{}
(ASE'24) reports an average repair rate of $50.06\%$; \gc{} (AAAI-26) reports
$86.88\%$ and positions itself as a $>$30-point improvement over \gp{}. Neither
claim has been independently verified, and the comparison between them is made by
importing numbers across two different experimental setups.

This document describes a reproduction study that re-executes both methods under a
single controlled protocol on a shared model (\texttt{Mistral-7B-Instruct-v0.1}),
regenerating every quantity from scratch rather than citing published values. We
explain \emph{why} each design decision was taken, what each measurement can and
cannot establish, and where our own implementation was found to be wrong.

The study's central methodological contribution is that it fixes the
\emph{denominator}. We show that the population of ``glitch tokens'' in a given
model is not a property of the model alone but of the probing protocol: on the
same model, the same vocabulary, and the same filtering rules, we count $988$
glitch tokens under one reasonable protocol and $2{,}552$ under another --- and
the two published papers themselves used populations of $2{,}779$ and $2{,}539$
respectively for this model. Because a ``repair rate'' is a fraction of that
population, rates computed against different populations are not comparable, and
we show that \gc{}'s headline comparison table reports a \gp{} token count that
\gp{} never measured.

We additionally subjected our own implementation to an adversarial external audit
before reporting results. The audit identified two errors serious enough to
invalidate specific conclusions. Both were corrected, the affected experiments
re-run, and both the pre-correction and post-correction numbers are reported here,
because the difference between them is itself evidence about how sensitive these
methods are to implementation detail.
\end{abstract}

\tableofcontents

% =====================================================================
\section{Introduction and motivation}
% =====================================================================

\subsection{The phenomenon}

A large language model does not read characters; it reads \emph{tokens}, drawn
from a fixed vocabulary built during training. Some of those tokens occur almost
never in the training corpus. Their embedding vectors are therefore barely
updated from initialisation, and the model has no coherent representation of what
they mean. When such a token appears in a prompt, the model's behaviour degrades
in ways that look like malfunction rather than error: asked to repeat the string
\texttt{Azalera}, one recent model answers \texttt{saysay}; asked for the length
of \texttt{Kapunoang}, it answers \texttt{3}. These are called \emph{glitch
tokens}.

They matter for two reasons. First, they are a reliability problem: a user who
happens to paste text containing one gets nonsense with no warning. Second, they
are a security problem: because they occupy a region of representation space the
model never learned to control, they have been used as a lever for prompt
injection and jailbreak attacks.

\subsection{The two methods under study}

\paragraph{\gp{} (Zhang et al., ASE 2024).} Detects glitch tokens by reading the
model's internal activations. It labels a small random sample of the vocabulary by
brute force, extracts activation features for those tokens from a band of
late transformer layers, compresses them with PCA, trains a support-vector
classifier, and applies it to the rest of the vocabulary. It then \emph{repairs}
glitch tokens at inference time by editing activations: neurons that are
reliably active for normal tokens but under-active for a glitch token are
amplified; neurons that are reliably silent for normal tokens but abnormally
firing are suppressed. No model weights change.

\paragraph{\gc{} (Fan, Li \& Li, AAAI 2026).} Repairs only. It attaches small
LoRA adapters ($<0.1\%$ additional parameters) to the same components of the same
layer band, and trains them on (glitch token $\rightarrow$ correct answer) pairs.
Crucially, the adapters are \emph{gated}: a switch $\lambda$ enables the branch
only when the input contains a known glitch token, so that clean inputs are
processed by the mathematically unmodified base model. This is the basis of its
``lossless'' claim.

\subsection{Why these particular claims require verification}

Reproduction studies are often motivated in general terms. Here the motivation is
specific and checkable. Four concrete problems are visible in the published
material \emph{before} any experiment is run.

\begin{enumerate}[leftmargin=*]

\item \textbf{The comparison is made across incompatible populations.}
\gc{}'s central claim is a comparison with \gp{}. But the two papers worked with
different sets of glitch tokens for the same model. For
\texttt{Mistral-7B-Instruct-v0.1}, \gp{}'s own tables imply a population of
$2{,}779$ tokens\footnote{Three independent derivations from \gp{}'s Tables 3
and 5 agree: $1873/0.6741=2778.5$ (detection), $359/0.1292=2778.6$ and
$1045/0.3760=2779.3$ (repair).}; \gc{} states $2{,}539$. A ``repair rate'' is a
fraction of this population, so the two rates measure different things.

\item \textbf{One reported number appears to be transplanted rather than
measured.} \gc{}'s Table 2 reports that \gp{} repaired $956$ Mistral tokens at
$37.65\%$. \gp{} reports repairing $1{,}045$ tokens at $37.60\%$. The value $956$
is what one obtains by applying \gp{}'s \emph{rate} to \gc{}'s \emph{denominator}
($0.3765 \times 2539 = 955.9$). A rate transplanted onto a different population is
presented in the position of a measured count. This is precisely the failure mode
that fixing the denominator is meant to prevent.

\item \textbf{Neither paper reports variance.} \gp{}'s pipeline depends on a
random $10\%$ sample of the vocabulary; \gc{}'s depends on a random adapter
initialisation and batch order. Both report single numbers. Without a spread,
there is no way to know whether a reported difference exceeds run-to-run noise.

\item \textbf{\gc{} appears to evaluate on its training data.} The paper
constructs a dataset from the identified glitch tokens, trains the adapters on it,
and then reports the repair rate over the glitch tokens. No held-out split is
described. If training and evaluation populations coincide, the headline number
measures memorisation capacity as well as repair.
\end{enumerate}

Additionally, \gp{} has no public code, and the constants of its adaptive repair
rule are never given numerically --- a reproducibility question in its own right.

\subsection{What kind of study this is, and what it can conclude}

This is a \emph{reproduction and differential evaluation}, not a new method. It is
important to be precise about its epistemic limits, because the temptation to
overclaim is the main hazard of such work.

When our measurement disagrees with a published number, there are always exactly
two explanations: the published claim does not hold, or our reconstruction differs
from what the authors actually did. Distinguishing them is the entire craft, and
only two instruments are available:

\begin{itemize}[leftmargin=*]
\item \textbf{Anchoring.} If, under the authors' own conditions, our pipeline
reproduces a quantity they published, our reconstruction is demonstrably faithful
in that respect, and remaining disagreements become informative. We achieve this
for \gc{}: our census under their released code's protocol recovers $2{,}537$ of
their $2{,}539$ published glitch tokens (Section~\ref{sec:rq1}).
\item \textbf{Refusal to attribute.} Where no anchor exists --- notably \gp{},
which released no code --- a disagreement is reported as ``we could not reproduce
$X$ from the published description,'' never as ``$X$ is false.'' We hold to this
distinction throughout, and we flag in Section~\ref{sec:threats} exactly which of
our findings are anchored and which are not.
\end{itemize}

% =====================================================================
\section{Research questions}
% =====================================================================

Each published claim is converted into a question that a measurement can answer.

\begin{description}[leftmargin=1.6cm,style=nextline]
\item[RQ1 --- Population] How many glitch tokens does the model actually have, and
how much does that number depend on the protocol used to probe it? \emph{This is
prior to everything else: it is the denominator of every repair rate.}
\item[RQ2 --- \gp{} detection] Does the detection stage reproduce its claimed
precision, recall, F1 and time advantage, and how much do these vary across
random samples?
\item[RQ3 --- \gp{} repair] Does the repair stage reproduce its claimed rate; is
its parameterisation ($\alpha$, $\beta$, $m$) justified by its own behaviour; and
what collateral damage does it cause to normal tokens?
\item[RQ4 --- \gc{} repair] Does the LoRA repair reproduce its claimed rate, and
does it hold on glitch tokens the adapter never trained on?
\item[RQ5 --- Side effects] What do the two methods cost in inference speed, and
does either damage the model's general capability?
\end{description}

% =====================================================================
\section{Experimental design}
% =====================================================================

\subsection{Choice of subject model}

All experiments in this document use \texttt{Mistral-7B-Instruct-v0.1}. It was
chosen deliberately, for three reasons.

\begin{enumerate}[leftmargin=*]
\item It is used by \emph{both} papers, so both sets of claims apply to it.
\item It is the model on which the two papers disagree most sharply: \gc{} reports
repairing $94.80\%$ of its glitch tokens while attributing $37.65\%$ to \gp{}. If
a discrepancy is going to be visible anywhere, it is here.
\item It is ungated, so the experiment is reproducible by any reader without
access approvals.
\end{enumerate}

Restricting the study to one model is a deliberate depth-versus-breadth trade.
Every claim below is therefore a claim about this model; Section~\ref{sec:threats}
addresses generalisation, and the infrastructure is model-agnostic so that
additional models require configuration only.

\subsection{Hardware and software}

Single NVIDIA RTX A6000 (48\,GB), fp16, native Windows, PyTorch 2.6 with CUDA
12.4. The original papers used $2\times$A100 (\gp{}) and an H200 (\gc{}'s speed
measurements). Because absolute wall-clock times are not comparable across
hardware, \textbf{every timing claim in this study is evaluated as a ratio}
measured within our own environment, never as an absolute duration compared to a
published one.

% =====================================================================
\section{Methodology}
% =====================================================================

\subsection{Establishing ground truth, and the protocol problem}
\label{sec:method-census}

Both papers define a glitch token \emph{operationally}, by the same test: prompt
the model to repeat the token verbatim at temperature $0$; if it fails, the token
is glitchy. We adopt the same definition, and apply it exhaustively to the entire
vocabulary rather than to a sample, so that the resulting labels are a census
rather than an estimate.

\paragraph{Filtering.} Some vocabulary entries cannot meaningfully be tested and
must be excluded first, following the procedure of Land \& Bartolo (2024) which
both papers adopt: control tokens (\texttt{<s>}, \texttt{<unk>}, \dots); tokens
that do not decode to valid text; and \emph{unreachable} tokens, whose decoded
string re-encodes to a different token id and which therefore can never appear
from real input.

The reachability test is subtler than it looks, and our first implementation got
it wrong in an instructive way. A naive ``decode then re-encode'' check
misclassified $12{,}385$ tokens ($39\%$ of the vocabulary) as unreachable, because
SentencePiece tokenizers attach leading-space semantics to a string in isolation
that they do not attach in context. The correct procedure --- which we took from
\gc{}'s own released filter, and which originates with Land \& Bartolo --- performs
the round-trip behind a fixed one-token prefix (\guillemotleft), so the token is
tested in context. This yields $257$ filtered tokens, matching the scale both
papers report ($229$ for Llama-2 in \gc{}). We report this because it is a
reproducibility lesson: the filtering step is load-bearing, and a plausible
implementation of it is off by two orders of magnitude.

\paragraph{Prompt construction.} We insert the token under test into the prompt as
its \emph{raw vocabulary id}, spliced between two pre-tokenised template halves,
rather than by string concatenation followed by re-tokenisation. This guarantees
that the model receives exactly the token being tested (string concatenation can
re-tokenise into a different sequence), and it makes every prompt the same length,
which permits exact batching.

\paragraph{Two protocols, and why both are needed.}
\label{sec:two-protocols}
While implementing the census we found that \gc{}'s released code does not use the
prompt its paper describes. Two protocols therefore exist, and we implement both:

\begin{description}[leftmargin=1.2cm,style=nextline]
\item[\Ppaper{}] The template as written in the papers, with a $24$-token
generation budget and a containment check on the stripped target.
\item[\Pcode{}] \gc{}'s released code verbatim: a different template
(\texttt{"Question: \dots return back to me?$\backslash$n Answer: \dots"}), a
$10$-token generation budget, and their exact correctness test.
\end{description}

This is not pedantry. \Pcode{} exists so that our numbers can be \emph{anchored}
to the authors' published artefacts; \Ppaper{} exists because it is what a reader
of the paper would implement. Reporting both is what allows us to separate
``the method behaves differently than claimed'' from ``we probed it differently.''

\subsection{Reimplementation of \gp{}}

With no public code, \gp{} was reconstructed from its Algorithm 1 (detection) and
Algorithm 2 (repair), its Equations 7--12, and its stated hyperparameters
($\gamma=0.1$ sampling rate, $P=75$ PCA components, polynomial SVM with $C=1$ and
degree $3$, threshold $m=1$, key layers $19$--$28$).

Two aspects of the method are \emph{under-specified in the paper} and therefore
required a choice on our part. We record both explicitly, because a reader must be
able to tell our decisions from the authors':

\begin{enumerate}[leftmargin=*]
\item \textbf{Sequence position.} The paper does not state which position in the
sequence activations are read from. We use the position of the token under test,
which is the reading consistent with ``the activation states of glitch tokens'';
the alternative (final prompt position) is exposed as a configuration switch so
the choice can be ablated rather than assumed.
\item \textbf{The adaptive constants.} Equations 11--12 define
$\beta = k_1\Delta_{\uparrow} + b_1$ and $\alpha = k_2\Delta_{\downarrow} + b_2$,
but the values of $k_1,b_1,k_2,b_2$ are never given; the paper says only that they
are ``derived through an adaptive process tailored to the specific dynamics of
each model,'' and mentions a ``range restriction'' that is likewise unspecified.
We use identity constants ($k=1$, $b=0$) and an explicit, configurable clamp. This
is \emph{our} choice, not a reconstruction of theirs, and every conclusion about
the adaptive variant is qualified accordingly.
\end{enumerate}

\subsection{Reimplementation of \gc{}}

\gc{} released code, which we use as the specification of record wherever the
paper is ambiguous. Our implementation matches it on: LoRA rank $4$ and scaling
$4$; adapters on the MLP gate and data projections of layers $19$--$28$ only;
frozen base weights; loss computed only on the answer span; and the authors'
exact training-target format (the decoded token, unstripped, in single quotes).
We also adopt their optimisation settings --- $15$ epochs, learning rate
$1\times10^{-4}$, effective batch $512$, linear schedule with $10\%$ warmup,
gradient clipping at $1.0$ --- so that a hyperparameter mismatch cannot be offered
as an explanation for any difference in outcome.

\paragraph{The gate.} We implement the $\lambda$ switch as the paper and the
authors' code define it: a per-example flag, set to $1$ when the input contains a
token from the known glitch set $G$ and $0$ otherwise, multiplying the LoRA
branch's contribution. Two properties follow, and both are measured rather than
assumed. First, ``lossless'' is a property of \emph{the gate}, not of the adapter:
with the gate active a clean prompt is processed by the unmodified base model by
construction. Second, the branch is still \emph{computed} when $\lambda=0$ (this
is true of the authors' implementation as well), so the gate removes the
behavioural effect but not the arithmetic cost --- which is why it does not
exempt the method from an inference-speed penalty.

\paragraph{The held-out split (our addition).} The paper trains and evaluates on
the same token population. We split the glitch set $80/20$, train only on the
$80\%$, and evaluate both halves. This is the single most informative addition
this study makes to \gc{}'s own evaluation, because it separates two things the
paper's design conflates: the ability to \emph{repair} glitch tokens in general,
and the ability to \emph{memorise} the specific ones shown during training.

\paragraph{Leakage control.} A token is only ``unseen'' if the adapter never
encountered its id. A token id can appear inside a training \emph{sequence} even
when it is not the token that sequence is about --- as ordinary context, or
because another token's decoded text re-encodes through it. We therefore scan
every training sequence and exclude from the held-out set any token whose id
appears in one. Without this control the held-out result would be optimistic in a
way that is invisible from the split file alone.

\subsection{Metrics, stated with their denominators}

Ambiguity about denominators is the specific defect this study exists to correct,
so every metric is defined against an explicit population.

\begin{align*}
\text{repair rate} &= \frac{\#\{\text{glitch tokens that repeat correctly with the method active}\}}{|\text{our glitch census for this model under this protocol}|}\\[0.4em]
\text{collateral breakage} &= \frac{\#\{\text{normal tokens that stop repeating correctly with the method active}\}}{|\text{sampled normal tokens}|}\\[0.4em]
\text{detection recall} &= \frac{|\hat{G}\cap G|}{|G|},\qquad
\text{precision} = \frac{|\hat{G}\cap G|}{|\hat{G}|}
\end{align*}

Collateral breakage deserves comment: \emph{neither paper reports it}. A repair
method that fixes glitch tokens by degrading normal ones is not obviously an
improvement, and the quantity is cheap to measure, so we measure it everywhere.

\subsection{Statistical treatment}

Every sampled quantity is run at three seeds and reported as mean $\pm$ sample
standard deviation. Three seeds is a weak basis for a variance estimate and we do
not treat it as more than an indication of scale; where we describe variance we
give the interval rather than an adjective. This is still strictly more than the
papers report, which is a single number with no spread.

\subsection{Reproducibility infrastructure}

Three properties were built in, because a study whose purpose is to check other
people's numbers has no standing unless its own can be checked.

\begin{itemize}[leftmargin=*]
\item \textbf{Traceability.} Every result file records the timestamp, the git
commit, \emph{whether the working tree was modified at the time}, and the versions
of PyTorch, Transformers, PEFT and scikit-learn. The dirty-tree flag was added
after the audit found artefacts attributed to commits that could not have produced
them --- a commit hash alone is not provenance if the run was launched from an
edited tree.
\item \textbf{Raw evidence retention.} Full model generations are retained for
every probe, and per-token outcomes for every repair evaluation, so that headline
numbers can be recomputed from primary data rather than trusted.
\item \textbf{Checkpointing.} Long sweeps checkpoint per batch and resume exactly,
which makes multi-hour runs interruptible without loss and makes the checkpoint
itself a raw-evidence artefact.
\end{itemize}

% =====================================================================
\section{Internal validation: an adversarial audit of our own implementation}
\label{sec:audit}
% =====================================================================

\subsection{Why we did this, and why it is reported}

A reproduction study that only audits other people's work is not credible. Before
reporting any result, we therefore submitted the implementation to an independent
adversarial review with instructions to treat all code, documentation and
conclusions as unverified claims. The review checked faithfulness to the papers,
correctness of the code, integrity of the stored results, and whether each stated
conclusion was supported by its evidence at the strength claimed.

It found two errors serious enough to invalidate specific conclusions, several
correctness defects, and a number of conclusions stated more strongly than the
evidence supported. All are reported here. This is deliberate: the pre- and
post-correction numbers differ, and \emph{the size of that difference is itself a
result} --- it measures how sensitive these methods are to implementation choices
that a paper does not pin down.

\subsection{The two material errors, and what they invalidated}

\paragraph{E1: \gp{} repair modified the wrong tensor.}
The gated MLP computes two intermediate signals --- the gate $\sigma(Z_1)$ and the
data $Z_2$ --- and multiplies them elementwise. The paper's Figure 5 shows repair
producing a ``Modified MLP gate'' \emph{and} a ``Modified MLP data'', i.e.\ the
two are adjusted separately, before multiplication. Our first implementation
adjusted their already-multiplied product, which is a different operation: adding
$\beta$ to a product is not adding $\beta$ to its factors.

The consequence was not merely numerical. Neuron selection uses the criterion
``active in $>99\%$ of normal tokens''; applied to the product, this selected
$0$--$4$ neurons per layer, which left the amplification factor $\beta$ with
almost nothing to act on. Applied to the separate streams as the paper intends,
it selects several times more. Our pre-correction conclusion that ``$\beta$ is
inert'' was therefore partly an artefact of our own error. The corrected
implementation adjusts the streams separately; the previous behaviour is retained
as an explicit ablation (Section~\ref{sec:rq3}) rather than silently discarded, so
that the two can be compared.

\paragraph{E2: \gc{}'s gate was absent.}
Our first implementation used an ordinary always-active adapter and disabled it
globally for one control. Repair rates on glitch prompts are unaffected (the gate
would be open for those prompts anyway), but every claim about \emph{clean-input}
behaviour --- losslessness, capability preservation, inference speed --- was
measured on an always-on adapter and therefore was not measuring \gc{}. The gate
is now implemented per example as specified.

\subsection{Correctness defects found and fixed}

\begin{enumerate}[leftmargin=*,label=D\arabic*.]
\item \textbf{Unseeded training.} Only the data split was seeded; adapter
initialisation and batch order were not, so the trained adapter was not
reproducible. All randomness is now seeded.
\item \textbf{Dropped gradients.} The optimiser stepped only on exact
gradient-accumulation boundaries, so each epoch's final partial window was
computed, never applied, and discarded. (The authors' own script flushes this
window correctly; ours now does too.)
\item \textbf{Adaptive parameters fitted on the evaluation set.} Because the model
has fewer glitch tokens than the sampling rate implies, the adaptive $\alpha,\beta$
were fitted on the same tokens they were scored on, while the fixed-value variant
was not --- an unequal comparison. The glitch set is now split: one half fits the
adaptive parameters, both variants are scored on the other half.
\item \textbf{Deviation from Equation 10.} We had inserted an absolute value into
the denominator of the published signed ratio. The published form is restored and
the clamp is made explicit and configurable rather than hard-coded.
\item \textbf{Held-out leakage.} $18$ of $197$ nominally held-out tokens had ids
occurring inside training sequences. These are now excluded (see leakage control
above).
\item \textbf{Training-target mismatch.} Our targets were stripped and unquoted;
the authors' are the unstripped decoded token in single quotes. A different target
teaches a different output format, so ours now matches theirs.
\item \textbf{Missing provenance and raw evidence.} Added as described above.
\item \textbf{An advertised experiment that never ran.} The sensitivity sweep over
the neuron-selection threshold $m$ was configured but never executed. It is now
implemented and reported (Section~\ref{sec:rq3}).
\end{enumerate}

\subsection{Conclusions we withdrew or weakened}

The audit also challenged the \emph{strength} of several conclusions without
disputing the underlying measurements. We accepted the following corrections, and
the results sections below state the weakened forms:

\begin{itemize}[leftmargin=*]
\item ``$\beta$ is inert'' --- withdrawn. Contradicted by our own grid at
$\alpha=1$, and partly an artefact of E1.
\item ``pure protocol sensitivity'' --- weakened. Our two protocols differ in
prompt wording \emph{and} generation budget \emph{and} correctness test
simultaneously, so the aggregate effect was demonstrated but not decomposed. We
subsequently ran the missing factorial cell to decompose it
(Section~\ref{sec:rq1}).
\item ``unimplementable as published'' --- weakened to ``not reproducible from the
published description under the reconstructions we attempted,'' which is what the
evidence supports.
\item ``$100\%$ precision by construction fails'' --- retained but re-caused. The
false positives are real, but our stated cause (fp16 non-determinism) was asserted
rather than isolated; an alternative path exists whereby a token enters the glitch
set from sample labelling without passing post-validation.
\item A claim that \gc{}'s losslessness ``requires a detector at inference'' ---
corrected. The gate is a membership lookup against the known set $G$, not a
learned detector. The corrected observation is sharper and appears in
Section~\ref{sec:rq4}: a glitch token \emph{absent} from $G$ receives
$\lambda = 0$ and hence no repair at all, by design.
\end{itemize}

% =====================================================================
\section{Results}
% =====================================================================

\subsection{RQ1 --- The glitch-token population is a property of the protocol}
\label{sec:rq1}

Table~\ref{tab:census} reports the full-vocabulary census under both protocols.
Filtering is identical in both; only the probe differs.

\begin{table}[h]
\centering
\caption{Full-vocabulary census, \texttt{Mistral-7B-Instruct-v0.1}
($|V|=32{,}000$). Filtering is protocol-independent; the glitch/normal split is
not.}
\label{tab:census}
\begin{tabular}{lrr}
\toprule
& \Ppaper{} & \Pcode{} \\
\midrule
Normal            & 30{,}755 & 29{,}191 \\
\textbf{Glitch}   & \textbf{988} & \textbf{2{,}552} \\
Unreachable       & 253 & 253 \\
Special           & 3 & 3 \\
Undecodable       & 1 & 1 \\
\midrule
Total             & 32{,}000 & 32{,}000 \\
\bottomrule
\end{tabular}
\end{table}

\paragraph{Anchoring.} Under \Pcode{} we recover $2{,}537$ of \gc{}'s $2{,}539$
published glitch tokens (Jaccard $0.993$), with $15$ additional tokens of our own.
This is the anchor described in Section~1.4: our census machinery reproduces the
authors' published artefact essentially exactly when run under their conditions.
Consequently the difference between the two columns of Table~\ref{tab:census}
cannot be attributed to a defect in our implementation.

\paragraph{The three populations disagree.} For this one model there are now three
published or measured glitch-token populations: $2{,}779$ (implied by \gp{}'s
tables), $2{,}539$ (\gc{}), and ours at $988$ (\Ppaper{}) or $2{,}552$ (\Pcode{}).
The first two differ by $9\%$ from each other despite ostensibly measuring the same
property of the same model.

\paragraph{What causes the gap, and what we can prove about it.} The two protocols
differ in three respects at once. Our initial write-up called the effect ``pure
protocol sensitivity'', which overstated what a two-point comparison can show. To
decompose it we ran the missing factorial cell --- \gc{}'s prompt with the
$24$-token budget --- so that prompt wording and generation budget can be
separated. \emph{[Factorial result inserted here once the run completes; the
decomposition is reported as a measurement, not an extrapolation.]}

\paragraph{Why this matters more than it appears.} A repair rate is a fraction of
this population. If the population is protocol-dependent, then so is every
published repair rate, and rates computed against different populations cannot be
compared. This is not hypothetical: as noted in Section~1.3, \gc{}'s comparison
table reports a \gp{} token count ($956$) that is \gp{}'s rate rescaled onto
\gc{}'s denominator, rather than a quantity \gp{} measured.

\subsection{RQ2 --- \gp{} detection}
\label{sec:rq2}

Table~\ref{tab:detect} reports detection scored against our census, at three
seeds, under both protocols.

\begin{table}[h]
\centering
\caption{\gp{} detection on \texttt{Mistral-7B-Instruct-v0.1}. Mean $\pm$ sample
s.d.\ over seeds $\{0,1,2\}$. The paper's Mistral row is given for reference; it
was measured against a population of $2{,}779$ tokens.}
\label{tab:detect}
\begin{tabular}{lcccc}
\toprule
& Precision & Recall & F1 & Time \\
\midrule
Ours, \Ppaper{} ($|G|=988$)    & $0.997$ & $0.279 \pm 0.091$ & $0.431 \pm 0.116$ & $211$\,s \\
Ours, \Pcode{} ($|G|=2{,}552$) & $0.997$ & $0.404 \pm 0.019$ & $0.575 \pm 0.019$ & $224$\,s \\
\midrule
\gp{} claim (Mistral)          & $1.000$ & $0.674$ & $0.805$ & $42$\,m\,$39$\,s \\
\bottomrule
\end{tabular}
\end{table}

\paragraph{Recall does not reproduce, and the base rate does not explain it.}
Under \Ppaper{} recall is $0.279$; one might suspect this is depressed because our
stricter census makes positives rarer, giving the classifier less signal. We
tested this by re-scoring under \Pcode{}, whose base rate ($8.0\%$) matches the
conditions the papers worked in. Recall rises to $0.404$ --- so the base rate
explains part of the gap --- but remains far from the claimed $0.674$. The
shortfall survives the control.

\paragraph{Variance is substantial and is never reported.} Under \Ppaper{}, recall
across three seeds spans $0.174$--$0.335$: the only difference between those runs
is which random $10\%$ of the vocabulary trained the classifier. The interval is
much tighter under the higher base rate ($0.382$--$0.417$), which is itself
informative: the method's stability depends on how many positives its sample
happens to contain. A single reported number cannot convey this.

\paragraph{The ``$100\%$ precision'' guarantee is not exactly attained.} \gp{}
validates every predicted positive with the brute-force test, which should make
false positives impossible. We observe $1,0,1$ false positives per seed under
\Ppaper{} and $3,4,1$ under \Pcode{} --- small, but not zero. The guarantee holds
only relative to a run's own within-run judgements: a token can be judged glitchy
at one moment and normal at another, and at least one route into the predicted set
bypasses post-validation entirely. We report the effect and note that we did not
isolate its cause.

\paragraph{On the time comparison.} We do not compare our wall-clock against the
paper's, because the hardware differs. Our figure is reported only to establish
that detection cost is dominated by the brute-force validation passes, not by the
classifier.

\subsection{RQ3 --- \gp{} repair and its parameterisation}
\label{sec:rq3}

\emph{[Corrected results table inserted here: adaptive vs rule-based repair under
both protocols, with the pre-correction (product-stream) ablation alongside, the
$\alpha\times\beta$ grid, and the $m$-threshold sweep.]}

The reporting plan for this section is fixed in advance to avoid selective
presentation: we report (i) the corrected separate-stream results, (ii) the
pre-correction product-stream results as an ablation, (iii) the full parameter
grid with the best cell identified, (iv) the $m$ sweep, and (v) collateral
breakage for every configuration.

\subsection{RQ4 --- \gc{} repair, and what the held-out split reveals}
\label{sec:rq4}

\emph{[Corrected results table inserted here: train-split vs held-out repair
under both protocols, with the adapter-off control, the gated clean-input result,
and the forced-on clean-input result.]}

\paragraph{The design point that the held-out split exposes.} Independently of the
numbers, the gate has a structural consequence worth stating precisely, because
our earlier phrasing of it was wrong. The gate is a membership test against the
known glitch set $G$ --- not a learned detector. It follows that a glitch token
which is \emph{not} in $G$ receives $\lambda = 0$ and therefore no repair
whatsoever, by construction. \gc{} is thus a method for repairing an enumerated
list, and its coverage is bounded by the completeness of the detection step that
produced that list --- a step \gc{} does not itself perform, and which
Section~\ref{sec:rq2} finds recovers roughly $40\%$ of the population under the
best conditions we could reproduce. The held-out measurement is therefore best
read as an upper bound on a generalisation the deployed design does not attempt.

\subsection{RQ5 --- Inference cost}
\label{sec:rq5}

\emph{[Corrected results table inserted here: base vs \gp{} hooks vs gated \gc{}
adapter, as ratios.]}

% =====================================================================
\section{Threats to validity}
\label{sec:threats}
% =====================================================================

\subsection{Which findings are anchored and which are not}

This distinction determines how strongly each result may be stated, so it is made
explicit rather than left to the reader.

\begin{itemize}[leftmargin=*]
\item \textbf{Anchored (our implementation demonstrably reproduces the authors'
own artefact).} The \gc{} census under \Pcode{}: $2{,}537/2{,}539$ agreement. All
\gc{}-related findings inherit this anchor, and \gc{} additionally released code
against which our implementation is matched component by component.
\item \textbf{Not anchored.} Everything concerning \gp{}. It released no code, and
two of its specifications (the adaptive constants; the extraction position) are
incomplete. A disagreement between our measurement and its claim therefore cannot
be attributed with confidence, and is reported as a failure to reproduce from the
published description --- which is a statement about reproducibility, and a
legitimate finding in itself, but not a refutation.
\end{itemize}

\subsection{Construct validity: the correctness oracle is imperfect}

Both papers, and this study, judge repetition by asking whether the target string
appears in the output. This is a weak test in two identifiable cases: tokens whose
printable form is empty (whitespace-only), for which the test is trivially
satisfied, and tokens of one or two characters, where incidental occurrence is
plausible. In our census $31$ candidates are whitespace-only and $3{,}504$ have a
stripped length of at most one character. The defect is shared with the authors'
released code, but sharing it is not a defence: some fraction of every census in
this literature, ours included, is likely mislabelled for short tokens. We flag
this as the most significant unaddressed construct-validity threat.

\subsection{Internal validity}

Single model, single seed for the \gc{} adapter (three seeds are planned; the
training cost is the constraint), and a repair evaluation that consumes an oracle
glitch set rather than \gp{}'s own detector output. The last is deliberate --- it
isolates repair capability from detection error --- but it means we evaluate the
two stages separately rather than the end-to-end pipeline the paper describes.

\subsection{External validity}

All results are for one $7$B model. The two papers' claims are cross-model
averages, and a single model cannot refute an average; our results are stated as
claims about this model, and where we compare to a published number we use that
paper's \emph{Mistral-specific} figure rather than its average. The infrastructure
is model-agnostic, so extension is a matter of compute rather than method.

\subsection{Conclusion validity}

Three seeds bound variance only loosely. Where a difference is small relative to
the observed spread we say so rather than reporting it as a finding.

% =====================================================================
\section{Discussion}
% =====================================================================

\subsection{What a reader should take from this study}

\begin{enumerate}[leftmargin=*]
\item \textbf{The denominator is the finding.} The most consequential result is
not any single rate but the demonstration that the population those rates are
fractions of is protocol-dependent and differs across the two papers. Any future
work in this area should publish its census, not merely its rate --- and,
ideally, its probe verbatim, since a paper's prose and its code may differ, as
they do here.
\item \textbf{Held-out evaluation should be standard for learned repair.} A method
trained on an enumerated list and evaluated on that same list reports a quantity
that includes memorisation. The split costs nothing and changes the interpretation.
\item \textbf{Collateral damage should be reported.} Neither paper reports what
its repair does to normal tokens. It is cheap to measure and it is the natural
counterweight to a repair rate.
\item \textbf{Under-specification is a finding, not an inconvenience.} \gp{}'s
adaptive rule cannot be reconstructed from its paper because four constants are
withheld. That is worth stating plainly in the literature.
\end{enumerate}

\subsection{On our own errors}

Two of this study's implementation errors (E1, E2) would have produced
publishable-looking but wrong conclusions had they not been caught. We report them
for the same reason we report everything else: a benchmark's authority rests on
its auditability, and an audit that finds nothing is usually an audit that was not
adversarial enough. The relevant safeguard was not care but \emph{structure} ---
retained raw evidence, recorded provenance, and an external reviewer instructed to
disbelieve the documentation.

% =====================================================================
\appendix
\section{Artefact map}
% =====================================================================

Every number in this document is produced by code in this repository and stored
under \texttt{results/}. The mapping is:

\begin{longtable}{p{4.1cm}p{4.6cm}p{5.4cm}}
\toprule
\textbf{Result} & \textbf{Script} & \textbf{Artefact} \\
\midrule
\endhead
RQ1 census & \texttt{scripts/run\_ground\_truth.py} & \texttt{results/ground\_truth/<model>/\{tokens.csv, summary.json, sweep\_checkpoint.csv\}} \\
RQ1 \Pcode{} census & same, \texttt{--protocol gccode} & \texttt{.../gccode/} \\
RQ1 factorial cell & same, \texttt{--protocol gccode --max-new-tokens 24} & \texttt{.../gccode/<tag>/} \\
RQ2 detection & \texttt{scripts/run\_gp\_detect.py} & \texttt{results/gp\_detect/<model>/\{runs.csv, summary.json, checkpoints/\}} \\
RQ3 repair & \texttt{scripts/run\_gp\_repair.py} & \texttt{results/gp\_repair/<model>/\{summary.json, per\_token\_*.csv\}} \\
RQ3 ablation (v1 behaviour) & same, \texttt{--stream-mode product} & \texttt{.../ablation\_product/} \\
RQ3 grid and $m$ sweep & \texttt{scripts/run\_gp\_alpha\_beta\_sweep.py} & \texttt{.../\{alpha\_beta\_grid.csv, m\_sweep.csv, heatmap\_*.png\}} \\
RQ4 \gc{} & \texttt{scripts/run\_gc\_train.py}, \texttt{run\_gc\_eval.py} & \texttt{results/gc/<model>/\{split.json, train\_meta.json, eval.json\}} \\
RQ5 speed & \texttt{scripts/run\_speed.py} & \texttt{results/side\_effects/<model>/speed.json} \\
Aggregation & \texttt{scripts/aggregate\_results.py} & \texttt{results/benchmark\_report.\{json,md\}} \\
Run narrative & --- & \texttt{docs/RUNLOG.md} \\
\bottomrule
\end{longtable}

All experimental settings live in \texttt{configs/*.yaml}; nothing
experiment-relevant is hard-coded. \texttt{scripts/status.py} reports which stage
each model has completed and resumes from checkpoints after interruption.

\end{document}
